\chapter{Tentatives échouées, perfectibles ou inachevées dans l'alimentation de la base}

Les réflexions sur les différents moyens d'alimenter la base de données ont surtout fonctionné par tâtonnements. Il était donc inévitable que ce travail mène aussi à des échecs ou à des idées qui, faute de temps, n'ont pas abouti. 

\section{Les limites de l'IA dans le traitement des factures}

Si les tests réalisés avec l'IA ont donné des résultats très satisfaisants, ils ont aussi mis en évidence les limites techniques des différents outils, même les plus performants. 

\vspace{1,5cm}
\emph{Gemini}\index{Gemini} reste limité sur certains aspects, c'est pourquoi il n'est pas en mesure d'effectuer à lui seul tous les traitements nécessaires à l'import des données dans la base. Comme nous l'avons évoqué dans le chapitre précédent, \emph{Gemini}\index{Gemini} n'est pas capable d'interroger les catalogues de la BnF. Dans l'espoir de lui faire remplir l'intégralité des champs de la table \enquote{ligne de facture}, nous lui avons demandé de relier les oeuvres présentes dans les factures aux notices correspondantes dans les catalogues de la BnF. L'objectif était que \emph{Gemini}\index{Gemini} remplisse la section \enquote{informations catalogue}, en renseignant les ARK des auteurs et des exemplaires, ainsi que leurs cotes dans des colonnes prévues à cet effet. À cette demande, l'outil tente de réaliser cette tâche et nous indique qu'il l'a faite, mais les résultats sont toujours très fantaisistes. Les cotes ne sont jamais les bonnes et les ARK ne renvoient à aucune page existante. On voit bien que \emph{Gemini}\index{Gemini} est capable de manipuler les informations qui lui sont données, mais ne sait pas en utiliser d'autres. 

Les limites de \emph{Gemini}\index{Gemini} ont aussi été constatées sur des traitements de données en masse. La possibilité d'accéder aux numérisations a permis d'envisager des traitements de données plus massifs, toujours pour alimenter la base. Dans les traitements précédents, il était toujours nécessaire, comme nous l'avons évoqué, de saisir manuellement un enregistrement \enquote{facture} pour pouvoir produire un identifiant \emph{Heurist} permettant de lier la facture à toutes les lignes qui la composent. Nous nous sommes donc demandés s'il était possible d'automatiser aussi la transcription des données pour remplir la table facture, puis de lier par l'intermédiaire de \emph{Gemini} les identifiants \emph{Heurist} des factures aux lignes de factures correspondantes. Cela servirait à automatiser le traitement successif de deux tables essentielles de notre base. Pour ce faire, nous avons tenté de procéder en deux prompts.
Le premier a consisté à demander à \emph{Gemini}\index{Gemini} de transcrire et structurer toutes les informations utiles au remplissage de la table \enquote{facture}. Pour cela, un fichier PDF de factures numérisées a été joint au prompt suivant :

\begin{quote}
    Voici un ensemble de factures numérisées et produites par le même libraire  

En ne prenant en compte que les factures (donc pas les relevés ni les correspondances), peux-tu en extraire certaines informations sous forme de tableur ? 

Structure l'informations dans les colonnes suivantes : date de la facture, prix total, frais, commentaire facture, dossier, librairie H-ID (laisser la colonne vide)

Dans la colonne "dossier", mettre la valeur suivante : RCOL-1(<À COMPLÉTER EN FONCTION DU DOSSIER>)
\end{quote}

Ce prompt a fonctionné. Malgré quelques corrections à effectuer, le fichier CSV produit contenait les informations demandées sur les trente-trois factures numérisées du libraire Georges Chrétien et a pu être importé dans la base. Nous avons donc fourni à \emph{Gemini}\index{Gemini} l'instruction servant à remplir \enquote{ligne de facture}. En plus du dossier numérisé, nous avons joint à ce prompt un export CSV, issu de la base, des dates et des identifiants \emph{Heurist} des trente-trois enregistrements \enquote{facture} correspondants. Le prompt contenait une instruction supplémentaire demandant à \emph{Gemini}\index{Gemini} d'analyser les dates des factures et de les lier aux identifiants \emph{Heurist}\index{Gemini} issus de l'export CSV correspondant à ces dates. \emph{Gemini}Gemini\index{Gemini} n'a pas du tout pu réaliser cette tâche, répondant simplement qu'il n'était qu'un modèle de langage et que cette demande dépassait ses compétences. Le test a ensuite été simplifié : le fichier CSV contenant les identifiants \emph{Heurist} a été retiré des pièces jointes, et les instructions qui allaient avec ont été supprimées. Cela visait à voir si Gemini parvenait à remplir la table \enquote{ligne de facture} à partir d'un dossier entier numérisé. Cette fois, \emph{Gemini}\index{Gemini} a tenté de réaliser la tâche, mais le résultat n'était pas satisfaisant, car toutes les lignes n'ont pas été transcrites et dans les rares cas où elles l'ont été, il y avait beaucoup trop d'erreurs. 

Plusieurs outils expérimentés n'ont pas donné de résultats convaincants. Tous nos tests sont détaillés dans une documentation spécifique, dont le lien est accessible dans les annexes\footnote{Voir \enquote{\nameref{doc IA}}}. Nous avons par exemple essayé l'outil \emph{Corpusense}\index{Corpusense}, fonctionnant sur l'utilisation de \emph{Mistral AI}\index{Mistral AI}, autre modèle de langage\footnote{\emph{Corpusense} a été expérimenté avec l'aide de Joseph Chazalon, enseignant-chercheur au laboratoire de recherche de l'EPITA et de Jonathan Perrinet, ingénieur dans ce même laboratoire.}. Nous ne nous sommes pas attardés sur cet outil, qui n'est pas idéal pour traiter les factures Rondel, mais il a produit certains résultats intéressants, comme indiqué dans la documentation. 

\section{Idées non retenues}

Dans l'espoir de faciliter l'exploitation du corpus par une solution \emph{Open Source}, une solution pour améliorer l'utilisation de l'outil \emph{NuExtract3} présenté précédemment a aussi été envisagée. Cette suggestion a été faite par Fantin Le Ber, chargé de projet de recherche au département des Arts du spectacle dans le cadre du projet TORNE-H, visant à intégrer l'intelligence artificielle dans les processus de traitement documentaire et scientifique des collections muséales\footcite{noauthor_torne-h_nodate}. Cette solution consistait à créer une application intégrant les traitements réalisés par \emph{NuExtract3}\index{NuExtract3}, grâce à Streamlit\index{Streamlit}, un outil \emph{Open Source} permettant de développer des applications grâce à \emph{Python}\footcite{noauthor_streamlit_nodate}. L'application envisagée prévoyait d'entrer une facture, de la traiter par \emph{NuExtract3} et d'en extraire directement un fichier CSV prêt à être importé dans la base. Cependant, du fait du manque de temps et des résultats peu prometteurs de cette solution d'IA par rapport à ceux de \index{Gemini}\index{Gemini} sur les factures Rondel, cette idée a rapidement été écartée. 

Il faut surtout retenir de ce dernier chapitre qu'il s'intéresse à la réflexion autour de ce qui a échoué ou n'a pas abouti dans ce travail. Ces tentatives sont essentielles à prendre en compte pour comprendre tout le reste de notre raisonnement, mais aussi pour avoir des pistes d'améliorations en vue d'une suite de ce projet. En effet, vu la vitesse à laquelle progresse l'IA dans la transcription et la structuration de données, ce qui n'est pas possible techniquement à ce stade le sera peut-être ultérieurement. 